\documentclass[11pt,a4paper]{article}
\usepackage[utf8]{inputenc}
\usepackage[T1]{fontenc}
\usepackage[spanish,es-minimal,es-noshorthands]{babel}
\AtBeginDocument{\renewcommand{\sin}{\operatorname{sen}}\renewcommand{\arcsin}{\operatorname{arcsen}}}
\usepackage{lmodern}
\usepackage[sfdefault,lining]{FiraSans}
\usepackage[a4paper,top=2.3cm,bottom=2.2cm,left=1.8cm,right=1.8cm,headheight=1.3cm,headsep=0.4cm,footskip=1cm]{geometry}
\usepackage{amsmath,amssymb}
\usepackage{icomma}
\usepackage{graphicx}
\usepackage[table]{xcolor}
\usepackage{tikz}
\usepackage{fontawesome5}
\usepackage{ccicons}
\usepackage{fancyhdr}
\usepackage{lastpage}
\usepackage{tasks}
\usepackage{needspace}
\usepackage[most]{tcolorbox}
\usepackage{array}
\usepackage[hidelinks]{hyperref}
\AtBeginDocument{\hypersetup{pdftitle={Ficha de ejercicios de la unidad 4: Ecuaciones de 1.er y 2.º grado. 2º ESO},pdfsubject={Matemáticas 2º ESO},pdfkeywords={ecuaciones, identidades, ecuaciones de primer grado, paréntesis, denominadores, problemas, ecuaciones de segundo grado, incompletas, fórmula, discriminante}}}

\definecolor{navy}{HTML}{1B2748}
\definecolor{accent}{HTML}{E8603A}
\definecolor{teal}{HTML}{1F9E90}
\definecolor{soft}{HTML}{EEF1F8}
\definecolor{muted}{HTML}{5D6685}

\setlength{\parindent}{0pt}
\setlength{\parskip}{3pt}
\settasks{label=\alph*), label-format=\bfseries\color{navy}, label-width=1.3em, item-indent=1.8em, column-sep=0.6em, after-item-skip=5pt, before-skip=3pt, after-skip=4pt}

% ---------- capçalera i peu ----------
\newcommand{\llicencia}{\href{https://creativecommons.org/licenses/by-nc-sa/4.0/deed.es}{\ccbyncsa}\ \ Licencia Creative Commons Atribución-NoComercial-CompartirIgual 4.0 Internacional (CC BY-NC-SA 4.0)}
\pagestyle{fancy}
\fancyhf{}
\fancyhead[L]{\raisebox{-0.35\height}{\includegraphics[height=0.95cm]{logo-gva.png}}}
\fancyhead[C]{\small\color{muted}Ficha de ejercicios · Unidad 4: Ecuaciones de 1.\textsuperscript{er} y 2.\textsuperscript{o} grado · 2\textsuperscript{o} ESO}
\fancyhead[R]{\raisebox{-0.35\height}{\includegraphics[height=0.95cm]{logo-vermellar.png}}}
\fancyfoot[L]{\scriptsize\color{muted}\llicencia}
\fancyfoot[R]{\small\color{muted}\thepage}
\renewcommand{\headrulewidth}{0.4pt}
\renewcommand{\headrule}{\color{muted!40}\hrule height \headrulewidth}
\fancypagestyle{primera}{\fancyhf{}\renewcommand{\headrulewidth}{0pt}
  \fancyfoot[L]{\scriptsize\color{muted}\llicencia}
  \fancyfoot[R]{\small\color{muted}\thepage}}

% ---------- nivells i exercicis ----------
\newcommand{\estrelles}[1]{{\color{accent}\ifcase#1\or\faStar\faStar[regular]\faStar[regular]\or\faStar\faStar\faStar[regular]\or\faStar\faStar\faStar\fi}}
\newcounter{exer}
\newcommand{\ex}[1]{\par\needspace{4\baselineskip}\medskip\stepcounter{exer}\noindent{\color{navy}\bfseries\theexer.}\ \ {\small\estrelles{#1}}\ \ }

\newcommand{\seccio}[2]{\par\needspace{10\baselineskip}\bigskip
  \begin{tcolorbox}[colback=navy,colframe=navy,coltext=white,arc=2mm,boxrule=0pt,left=3mm,right=3mm,top=1.2mm,bottom=1.2mm]
  {\large\bfseries #1\quad #2}
  \end{tcolorbox}\vspace{-2pt}}

\newcolumntype{C}[1]{>{\centering\arraybackslash}p{#1}}
% quadre de recordatori (només a la fitxa d'exercicis)
\newcommand{\recorda}[1]{\par\smallskip\begin{tcolorbox}[colback=accent!7,colframe=accent!55,arc=2mm,boxrule=0.5pt,left=2.5mm,right=2.5mm,top=1mm,bottom=1mm]\small{\color{accent}\faLightbulb[regular]}\ \textbf{Recuerda:}\ #1\end{tcolorbox}}
\newcommand{\res}[1]{#1}\newcommand{\rest}[1]{#1}

\def\fmneed{12}
% =====================================================================
% Capçalera FentMates de les fitxes: un color per a cada curs.
% Els colors són els de la pàgina inicial (LEVELS a assets/js/materials.js);
% el logo ve de _eines/logo/pdf/ (python3 _eines/logo/genera.py).
%
% Ús, al final del preàmbul de la unitat:
%   % =====================================================================
% Capçalera FentMates de les fitxes: un color per a cada curs.
% Els colors són els de la pàgina inicial (LEVELS a assets/js/materials.js);
% el logo ve de _eines/logo/pdf/ (python3 _eines/logo/genera.py).
%
% Ús, al final del preàmbul de la unitat:
%   % =====================================================================
% Capçalera FentMates de les fitxes: un color per a cada curs.
% Els colors són els de la pàgina inicial (LEVELS a assets/js/materials.js);
% el logo ve de _eines/logo/pdf/ (python3 _eines/logo/genera.py).
%
% Ús, al final del preàmbul de la unitat:
%   \input{../comu/fentmates.tex}
%   \fmcurs{1eso}
% i, al començament del document:
%   \capcalera{Matemàtiques}{1r ESO}{Fitxa d'exercicis · Unitat 2}{Divisibilitat}
% Opcional abans de \input: \def\fmneed{14} (espai mínim abans de cada secció).
% =====================================================================
\makeatletter
\def\fm@logodir{../../logo/pdf}
\providecommand{\fmneed}{10}
\def\fm@cursnom{}

% clau del curs, color i color del text damunt del color
\newcommand{\fm@curs}[3]{\@namedef{fm@color@#1}{#2}\@namedef{fm@text@#1}{#3}}
\fm@curs{1eso}{E6B84A}{navy}
\fm@curs{2eso}{3B82F6}{white}
\fm@curs{3eso}{FF7A50}{white}
\fm@curs{4eso-a}{A78BFA}{white}
\fm@curs{4eso-b}{6D28D9}{white}
\fm@curs{1batx-mi}{9B1C1C}{white}
\fm@curs{1batx-ccss}{EF6B6B}{white}
\fm@curs{2batx-mii}{1E6B47}{white}
\fm@curs{2batx-ccss}{2BB3A3}{white}

\newcommand{\fmcurs}[1]{%
  \@ifundefined{fm@color@#1}{\PackageError{fentmates}{Curs desconegut: #1}{}}{}%
  \edef\fm@tmp{\noexpand\definecolor{curs}{HTML}{\@nameuse{fm@color@#1}}%
               \noexpand\colorlet{curstext}{\@nameuse{fm@text@#1}}}\fm@tmp
  \colorlet{cursfosc}{curs!72!black}%
  \def\fm@logo{\fm@logodir/vertical-#1.pdf}%
  \def\fm@marca{\fm@logodir/marca-#1.pdf}}

% ---------- capçalera de la primera pàgina ----------
% #1 matèria (petit), #2 curs, #3 tipus de document i unitat, #4 títol
\newcommand{\capcalera}[4]{%
  \gdef\fm@cursnom{#2}%
  \par\noindent
  \begin{tcolorbox}[colback=curs,colframe=curs,boxrule=0pt,arc=2.2mm,width=3.5cm,height=2.15cm,
      valign=center,left=3mm,right=2mm,top=1mm,bottom=1mm,nobeforeafter,box align=center]
    \color{curstext}{\footnotesize #1}\par\vspace{1pt}{\fontsize{19}{21}\selectfont\bfseries #2\par}%
  \end{tcolorbox}\hspace{5mm}%
  \begin{minipage}[c]{\dimexpr\linewidth-3.5cm-5mm-2.9cm\relax}
    {\color{cursfosc}\small\bfseries\MakeUppercase{#3}}\par\vspace{2pt}
    {\color{navy}\fontsize{20}{23}\selectfont\bfseries #4\par}
  \end{minipage}\hfill
  \begin{minipage}[c]{2.6cm}\raggedleft\includegraphics[width=2.6cm]{\fm@logo}\end{minipage}\par
  \vspace{2pt}}

% ---------- barres de secció del color del curs ----------
\renewcommand{\seccio}[2]{\par\needspace{\fmneed\baselineskip}\bigskip
  \begin{tcolorbox}[colback=curs,colframe=curs,coltext=curstext,arc=2mm,boxrule=0pt,left=3mm,right=3mm,top=1.2mm,bottom=1.2mm]
  {\large\bfseries #1\quad #2}
  \end{tcolorbox}\vspace{-2pt}}

% ---------- estrelles de dificultat del color del curs (a Batxillerat no se'n posen) ----------
\@ifundefined{estrelles}{}{%
  \renewcommand{\estrelles}[1]{{\color{curs}\ifcase#1\or\faStar\faStar[regular]\faStar[regular]\or\faStar\faStar\faStar[regular]\or\faStar\faStar\faStar\fi}}}

% ---------- capçalera i peu de les pàgines ----------
% primera pàgina: logos institucionals al peu
\fancypagestyle{primera}{\fancyhf{}\renewcommand{\headrulewidth}{0pt}
  \fancyfoot[L]{\raisebox{-0.32\height}{\includegraphics[height=0.78cm]{logo-gva.png}}\hspace{5mm}%
    \raisebox{-0.32\height}{\includegraphics[height=0.78cm]{logo-vermellar.png}}}
  \fancyfoot[R]{\scriptsize\color{muted}\href{https://creativecommons.org/licenses/by-nc-sa/4.0/deed.ca}{\ccbyncsa}\ \ CC BY-NC-SA 4.0 · fentmates.org\qquad\small\thepage}}
% la resta: el símbol a l'esquerra i el curs a la dreta, amb el títol al centre com abans
\fancyhead[L]{\raisebox{-0.3\height}{\includegraphics[height=0.62cm]{\fm@marca}}}
\fancyhead[R]{\small\bfseries\color{cursfosc}\fm@cursnom}
\renewcommand{\headrule}{\color{curs!45}\hrule height \headrulewidth}
\makeatother

%   \fmcurs{1eso}
% i, al començament del document:
%   \capcalera{Matemàtiques}{1r ESO}{Fitxa d'exercicis · Unitat 2}{Divisibilitat}
% Opcional abans de \input: \def\fmneed{14} (espai mínim abans de cada secció).
% =====================================================================
\makeatletter
\def\fm@logodir{../../logo/pdf}
\providecommand{\fmneed}{10}
\def\fm@cursnom{}

% clau del curs, color i color del text damunt del color
\newcommand{\fm@curs}[3]{\@namedef{fm@color@#1}{#2}\@namedef{fm@text@#1}{#3}}
\fm@curs{1eso}{E6B84A}{navy}
\fm@curs{2eso}{3B82F6}{white}
\fm@curs{3eso}{FF7A50}{white}
\fm@curs{4eso-a}{A78BFA}{white}
\fm@curs{4eso-b}{6D28D9}{white}
\fm@curs{1batx-mi}{9B1C1C}{white}
\fm@curs{1batx-ccss}{EF6B6B}{white}
\fm@curs{2batx-mii}{1E6B47}{white}
\fm@curs{2batx-ccss}{2BB3A3}{white}

\newcommand{\fmcurs}[1]{%
  \@ifundefined{fm@color@#1}{\PackageError{fentmates}{Curs desconegut: #1}{}}{}%
  \edef\fm@tmp{\noexpand\definecolor{curs}{HTML}{\@nameuse{fm@color@#1}}%
               \noexpand\colorlet{curstext}{\@nameuse{fm@text@#1}}}\fm@tmp
  \colorlet{cursfosc}{curs!72!black}%
  \def\fm@logo{\fm@logodir/vertical-#1.pdf}%
  \def\fm@marca{\fm@logodir/marca-#1.pdf}}

% ---------- capçalera de la primera pàgina ----------
% #1 matèria (petit), #2 curs, #3 tipus de document i unitat, #4 títol
\newcommand{\capcalera}[4]{%
  \gdef\fm@cursnom{#2}%
  \par\noindent
  \begin{tcolorbox}[colback=curs,colframe=curs,boxrule=0pt,arc=2.2mm,width=3.5cm,height=2.15cm,
      valign=center,left=3mm,right=2mm,top=1mm,bottom=1mm,nobeforeafter,box align=center]
    \color{curstext}{\footnotesize #1}\par\vspace{1pt}{\fontsize{19}{21}\selectfont\bfseries #2\par}%
  \end{tcolorbox}\hspace{5mm}%
  \begin{minipage}[c]{\dimexpr\linewidth-3.5cm-5mm-2.9cm\relax}
    {\color{cursfosc}\small\bfseries\MakeUppercase{#3}}\par\vspace{2pt}
    {\color{navy}\fontsize{20}{23}\selectfont\bfseries #4\par}
  \end{minipage}\hfill
  \begin{minipage}[c]{2.6cm}\raggedleft\includegraphics[width=2.6cm]{\fm@logo}\end{minipage}\par
  \vspace{2pt}}

% ---------- barres de secció del color del curs ----------
\renewcommand{\seccio}[2]{\par\needspace{\fmneed\baselineskip}\bigskip
  \begin{tcolorbox}[colback=curs,colframe=curs,coltext=curstext,arc=2mm,boxrule=0pt,left=3mm,right=3mm,top=1.2mm,bottom=1.2mm]
  {\large\bfseries #1\quad #2}
  \end{tcolorbox}\vspace{-2pt}}

% ---------- estrelles de dificultat del color del curs (a Batxillerat no se'n posen) ----------
\@ifundefined{estrelles}{}{%
  \renewcommand{\estrelles}[1]{{\color{curs}\ifcase#1\or\faStar\faStar[regular]\faStar[regular]\or\faStar\faStar\faStar[regular]\or\faStar\faStar\faStar\fi}}}

% ---------- capçalera i peu de les pàgines ----------
% primera pàgina: logos institucionals al peu
\fancypagestyle{primera}{\fancyhf{}\renewcommand{\headrulewidth}{0pt}
  \fancyfoot[L]{\raisebox{-0.32\height}{\includegraphics[height=0.78cm]{logo-gva.png}}\hspace{5mm}%
    \raisebox{-0.32\height}{\includegraphics[height=0.78cm]{logo-vermellar.png}}}
  \fancyfoot[R]{\scriptsize\color{muted}\href{https://creativecommons.org/licenses/by-nc-sa/4.0/deed.ca}{\ccbyncsa}\ \ CC BY-NC-SA 4.0 · fentmates.org\qquad\small\thepage}}
% la resta: el símbol a l'esquerra i el curs a la dreta, amb el títol al centre com abans
\fancyhead[L]{\raisebox{-0.3\height}{\includegraphics[height=0.62cm]{\fm@marca}}}
\fancyhead[R]{\small\bfseries\color{cursfosc}\fm@cursnom}
\renewcommand{\headrule}{\color{curs!45}\hrule height \headrulewidth}
\makeatother

%   \fmcurs{1eso}
% i, al començament del document:
%   \capcalera{Matemàtiques}{1r ESO}{Fitxa d'exercicis · Unitat 2}{Divisibilitat}
% Opcional abans de \input: \def\fmneed{14} (espai mínim abans de cada secció).
% =====================================================================
\makeatletter
\def\fm@logodir{../../logo/pdf}
\providecommand{\fmneed}{10}
\def\fm@cursnom{}

% clau del curs, color i color del text damunt del color
\newcommand{\fm@curs}[3]{\@namedef{fm@color@#1}{#2}\@namedef{fm@text@#1}{#3}}
\fm@curs{1eso}{E6B84A}{navy}
\fm@curs{2eso}{3B82F6}{white}
\fm@curs{3eso}{FF7A50}{white}
\fm@curs{4eso-a}{A78BFA}{white}
\fm@curs{4eso-b}{6D28D9}{white}
\fm@curs{1batx-mi}{9B1C1C}{white}
\fm@curs{1batx-ccss}{EF6B6B}{white}
\fm@curs{2batx-mii}{1E6B47}{white}
\fm@curs{2batx-ccss}{2BB3A3}{white}

\newcommand{\fmcurs}[1]{%
  \@ifundefined{fm@color@#1}{\PackageError{fentmates}{Curs desconegut: #1}{}}{}%
  \edef\fm@tmp{\noexpand\definecolor{curs}{HTML}{\@nameuse{fm@color@#1}}%
               \noexpand\colorlet{curstext}{\@nameuse{fm@text@#1}}}\fm@tmp
  \colorlet{cursfosc}{curs!72!black}%
  \def\fm@logo{\fm@logodir/vertical-#1.pdf}%
  \def\fm@marca{\fm@logodir/marca-#1.pdf}}

% ---------- capçalera de la primera pàgina ----------
% #1 matèria (petit), #2 curs, #3 tipus de document i unitat, #4 títol
\newcommand{\capcalera}[4]{%
  \gdef\fm@cursnom{#2}%
  \par\noindent
  \begin{tcolorbox}[colback=curs,colframe=curs,boxrule=0pt,arc=2.2mm,width=3.5cm,height=2.15cm,
      valign=center,left=3mm,right=2mm,top=1mm,bottom=1mm,nobeforeafter,box align=center]
    \color{curstext}{\footnotesize #1}\par\vspace{1pt}{\fontsize{19}{21}\selectfont\bfseries #2\par}%
  \end{tcolorbox}\hspace{5mm}%
  \begin{minipage}[c]{\dimexpr\linewidth-3.5cm-5mm-2.9cm\relax}
    {\color{cursfosc}\small\bfseries\MakeUppercase{#3}}\par\vspace{2pt}
    {\color{navy}\fontsize{20}{23}\selectfont\bfseries #4\par}
  \end{minipage}\hfill
  \begin{minipage}[c]{2.6cm}\raggedleft\includegraphics[width=2.6cm]{\fm@logo}\end{minipage}\par
  \vspace{2pt}}

% ---------- barres de secció del color del curs ----------
\renewcommand{\seccio}[2]{\par\needspace{\fmneed\baselineskip}\bigskip
  \begin{tcolorbox}[colback=curs,colframe=curs,coltext=curstext,arc=2mm,boxrule=0pt,left=3mm,right=3mm,top=1.2mm,bottom=1.2mm]
  {\large\bfseries #1\quad #2}
  \end{tcolorbox}\vspace{-2pt}}

% ---------- estrelles de dificultat del color del curs (a Batxillerat no se'n posen) ----------
\@ifundefined{estrelles}{}{%
  \renewcommand{\estrelles}[1]{{\color{curs}\ifcase#1\or\faStar\faStar[regular]\faStar[regular]\or\faStar\faStar\faStar[regular]\or\faStar\faStar\faStar\fi}}}

% ---------- capçalera i peu de les pàgines ----------
% primera pàgina: logos institucionals al peu
\fancypagestyle{primera}{\fancyhf{}\renewcommand{\headrulewidth}{0pt}
  \fancyfoot[L]{\raisebox{-0.32\height}{\includegraphics[height=0.78cm]{logo-gva.png}}\hspace{5mm}%
    \raisebox{-0.32\height}{\includegraphics[height=0.78cm]{logo-vermellar.png}}}
  \fancyfoot[R]{\scriptsize\color{muted}\href{https://creativecommons.org/licenses/by-nc-sa/4.0/deed.ca}{\ccbyncsa}\ \ CC BY-NC-SA 4.0 · fentmates.org\qquad\small\thepage}}
% la resta: el símbol a l'esquerra i el curs a la dreta, amb el títol al centre com abans
\fancyhead[L]{\raisebox{-0.3\height}{\includegraphics[height=0.62cm]{\fm@marca}}}
\fancyhead[R]{\small\bfseries\color{cursfosc}\fm@cursnom}
\renewcommand{\headrule}{\color{curs!45}\hrule height \headrulewidth}
\makeatother

\fmcurs{2eso}

\begin{document}
\thispagestyle{primera}

\capcalera{Matemáticas}{2\textsuperscript{o} ESO}{Ficha de ejercicios · Unidad 4}{Ecuaciones de 1.\textsuperscript{er} y 2.\textsuperscript{o} grado}
\vspace{10pt}
Nombre y apellidos: \rule{6.6cm}{0.4pt}\hfill Grupo: \rule{1.6cm}{0.4pt}\hfill Fecha: \rule{2.2cm}{0.4pt}
\vspace{8pt}

\begin{tcolorbox}[colback=soft,colframe=soft,arc=2mm,left=3mm,right=3mm,top=1.5mm,bottom=1.5mm]
\small Ejercicios graduados por niveles de dificultad:\quad \estrelles{1}\ básico \qquad \estrelles{2}\ medio \qquad \estrelles{3}\ avanzado.\\[2pt] Resuélvelos en el cuaderno mostrando todos los pasos y comprueba las soluciones siempre que puedas. Haz los cálculos sin calculadora. Las secciones siguen el orden de los apuntes de la unidad (4.1--4.9).
\end{tcolorbox}

% =====================================================
\seccio{1}{Ecuaciones: elementos y soluciones}

\ex{1} Indica qué igualdades son identidades y cuáles son ecuaciones:\begin{tasks}(2)
  \task $3(x+2)=3x+6$
  \task $3x+2=8$
  \task $(x+1)^{2}=x^{2}+2x+1$
  \task $x^{2}-1=0$
  \task $5x-2x=3x$
  \task $5x-2=3x$
  \task $(x-3)(x+3)=x^{2}-9$
  \task $2(x-1)=2x-1$
\end{tasks}\recorda{una \textbf{identidad} se cumple para cualquier valor de $x$; una \textbf{ecuación}, solo para algunos valores (o para ninguno).}

\ex{1} Comprueba si el valor indicado es solución de cada ecuación:\begin{tasks}(2)
  \task $3x-5=x+1$, \ $x=3$
  \task $x^{2}-4x=5$, \ $x=-1$
  \task $2(x-1)=x+4$, \ $x=5$
  \task $\dfrac{x}{2}+1=x-3$, \ $x=8$
  \task $x^{2}+x=6$, \ $x=-3$
  \task $7-x=2x+1$, \ $x=3$
\end{tasks}

\ex{1} En la ecuación $4x^2+x-3=2x+5$, indica el primer y el segundo miembro, los términos, la incógnita y el grado.

\ex{2} Indica el grado de cada ecuación una vez reducida:\begin{tasks}(2)
  \task $x(x+4)=x^{2}-8$
  \task $(x-2)^{2}=x^{2}+4$
  \task $x^{2}+3x=2x^{2}-1$
  \task $(x+1)(x-1)=3x$
\end{tasks}\recorda{antes de decir el grado, opera y reduce: los términos en $x^2$ se pueden cancelar.}

\ex{2} Escribe una ecuación de primer grado que tenga la solución $x=-2$ y una ecuación de segundo grado que tenga las soluciones $x=3$ y $x=-3$.

\ex{3} Halla el valor de $k$ en cada caso:\begin{tasks}(1)
  \task $x=3$ es solución de $kx-4=2x+k$.
  \task $x=-1$ es solución de $x^2+kx=k+5$.
\end{tasks}


% =====================================================
\seccio{2}{Ecuaciones de primer grado}

\ex{1} Resuelve mentalmente o en un solo paso:\begin{tasks}(4)
  \task $x+7=12$
  \task $x-5=-3$
  \task $3x=-21$
  \task $-4x=20$
  \task $\dfrac{x}{5}=3$
  \task $-x=6$
  \task $2x=7$
  \task $\dfrac{x}{4}=-2$
\end{tasks}\recorda{lo que suma pasa restando y lo que resta pasa sumando; lo que multiplica a todo un miembro pasa dividiendo (con su signo) y lo que divide pasa multiplicando.}

\ex{1} Resuelve:\begin{tasks}(3)
  \task $5x-4=11$
  \task $-2x+9=1$
  \task $3x+8=2$
  \task $7-x=10$
  \task $4x-3=-15$
  \task $6+5x=-4$
\end{tasks}

\ex{1} Resuelve:\begin{tasks}(3)
  \task $4x+3=2x+11$
  \task $7-3x=x-13$
  \task $6x-5=9x+7$
  \task $2x-1=5x-1$
  \task $8x+2=3x-13$
  \task $10-x=4x$
\end{tasks}

\ex{2} Reduce cada miembro y resuelve:\begin{tasks}(2)
  \task $5x-3+2x=4x+9$
  \task $8-3x+x=2x-4-x$
  \task $2x+7-5x=x-9$
  \task $x-4+3x=6x+2-x$
\end{tasks}

\ex{2} Resuelve (las soluciones son fracciones: simplifícalas):\begin{tasks}(4)
  \task $6x+1=2x+7$
  \task $5x-2=2x+3$
  \task $4-7x=x$
  \task $9x+4=3x-6$
\end{tasks}

\ex{3} Resuelve y comprueba la solución:\begin{tasks}(2)
  \task $3x-8+5x-1=4x+7-2x+2$
  \task $12-5x+3x=2x-9+x+1$
  \task $7x-2-x+5=4x+11-3x$
  \task $-x+6-4x=2-3x+x-10$
\end{tasks}

\ex{3} Aitana ha resuelto la ecuación $5x-3=2x+9$ así:\quad $5x+2x=9-3\ \Rightarrow\ 7x=6\ \Rightarrow\ x=\tfrac67$. \ ¿Qué errores ha cometido? Resuélvela correctamente y compruébala.


% =====================================================
\seccio{3}{Ecuaciones con paréntesis}

\ex{1} Resuelve quitando primero los paréntesis:\begin{tasks}(2)
  \task $3(x-2)=12$
  \task $5(x+1)=2x-7$
  \task $2(x-4)=x+1$
  \task $4(2x-1)=3x+6$
\end{tasks}

\ex{1} Resuelve:\begin{tasks}(4)
  \task $4-(x+3)=2x-5$
  \task $x-(5-x)=3$
  \task $7-(2x-1)=x+2$
  \task $-(x-6)=2x$
\end{tasks}\recorda{un signo $-$ delante de un paréntesis cambia el signo de \textbf{todos} los términos de dentro: $-(2x-1)=-2x+1$.}

\ex{2} Resuelve:\begin{tasks}(2)
  \task $3(2x-1)-2(x+4)=5x-13$
  \task $5-2(3-x)=4(x+1)-11$
  \task $4(x-3)-(x-5)=2(x+1)$
  \task $2(x+3)-3(x-1)=4x-6$
\end{tasks}

\ex{2} Resuelve y di qué ecuaciones no tienen solución y cuáles tienen infinitas:\begin{tasks}(2)
  \task $3(x+2)-x=2(x+3)$
  \task $5x-2(x-1)=3x+4$
  \task $6x-2(3x+1)=-2$
  \task $7x-3(2x-1)=2x+1$
\end{tasks}

\ex{2} Resuelve usando las identidades notables (verás que los términos en $x^2$ se cancelan):\begin{tasks}(3)
  \task $(x+1)^{2}-(x-1)^{2}=12$
  \task $x(x+3)-x^{2}=2x+5$
  \task $(x+2)(x-2)=x(x+1)-7$
\end{tasks}

\ex{3} Resuelve (empieza por los paréntesis de dentro):\begin{tasks}(2)
  \task $x-[2x-(3-x)]=7$
  \task $2[3-2(x+1)]=5-3(2x-1)-x$
\end{tasks}

\ex{3} Halla el valor de $a$ para que la ecuación $2(x+a)=x+9$ tenga la solución $x=3$.


% =====================================================
\seccio{4}{Ecuaciones con denominadores}

\ex{1} Resuelve:\begin{tasks}(4)
  \task $\dfrac{x}{3}=4$
  \task $\dfrac{x}{2}+\dfrac{x}{4}=6$
  \task $x-\dfrac{x}{3}=10$
  \task $\dfrac{2x}{5}-1=3$
\end{tasks}

\ex{1} Resuelve:\begin{tasks}(4)
  \task $\dfrac{x+2}{3}=4$
  \task $\dfrac{x-1}{2}=\dfrac{x+3}{4}$
  \task $\dfrac{3x-1}{4}=2$
  \task $\dfrac{x}{6}+1=\dfrac{x}{4}$
\end{tasks}\recorda{multiplica \textbf{todos} los términos por el m.c.m. de los denominadores (también los que no tienen denominador) y pon los numeradores entre paréntesis.}

\ex{2} Resuelve:\begin{tasks}(2)
  \task $\dfrac{x+1}{2}-\dfrac{x-1}{3}=2$
  \task $\dfrac{x-2}{4}-\dfrac{x-3}{6}=\dfrac{1}{2}$
  \task $\dfrac{2x+1}{3}-\dfrac{x-2}{2}=3$
  \task $x-\dfrac{x+2}{5}=\dfrac{x+6}{2}$
\end{tasks}

\ex{2} Resuelve:\begin{tasks}(2)
  \task $\dfrac{x-1}{3}+\dfrac{x+2}{6}=x-2$
  \task $3-\dfrac{x-1}{4}=\dfrac{x+5}{2}$
\end{tasks}

\ex{3} Resuelve e interpreta el resultado:\begin{tasks}(1)
  \task $\dfrac{3(x-2)}{4}-\dfrac{2(x+1)}{3}=\dfrac{x-10}{6}$
  \task $\dfrac{x}{2}-\dfrac{x-1}{3}-\dfrac{x+1}{4}=\dfrac{1-x}{12}$
  \task $\dfrac{x+1}{2}-\dfrac{x+3}{4}=\dfrac{x}{4}-1$
\end{tasks}

\ex{3} Un alumno ha resuelto la ecuación $\dfrac{x-2}{3}-\dfrac{x+1}{2}=1$ así:\quad $2x-4-3x+3=6\ \Rightarrow\ -x=7\ \Rightarrow\ x=-7$. \ ¿Dónde se ha equivocado? Resuélvela correctamente.


% =====================================================
\seccio{5}{Problemas con ecuaciones de primer grado}

\ex{1} La suma de dos números consecutivos es 75. ¿Cuáles son?\recorda{elige la incógnita y escribe qué representa, plantea la ecuación, resuélvela y comprueba que la solución tiene sentido en el problema.}

\ex{1} Si a un número le sumamos su triple, obtenemos 56. ¿Cuál es el número?

\ex{1} Ana tiene 12 años y su padre, 40. ¿Dentro de cuántos años la edad del padre será el doble que la de Ana?

\ex{2} El largo de un rectángulo es 3 m más que el doble del ancho y el perímetro mide 54 m. ¿Qué dimensiones tiene?

\ex{2} En un corral hay gallinas y conejos. En total hay 40 cabezas y 110 patas. ¿Cuántos animales hay de cada tipo?

\ex{2} En una hucha hay 30 monedas, unas de 1 € y otras de 2 €, que suman 46 €. ¿Cuántas monedas hay de cada tipo?

\ex{2} La edad de una madre es el triple de la de su hijo, y dentro de 12 años solo será el doble. ¿Cuántos años tienen ahora?

\ex{2} Después de un descuento del 15 \%, un abrigo cuesta 68 €. ¿Cuánto costaba antes del descuento?

\ex{3} Una persona gasta la tercera parte de su dinero en ropa y la cuarta parte en libros, y aún le quedan 25 €. ¿Cuánto dinero tenía?

\ex{3} Repartimos 150 € entre tres hermanos: el segundo recibe el doble que el primero, y el tercero, 30 € más que el segundo. ¿Cuánto recibe cada uno?

\ex{3} Queremos 20 kg de una mezcla de café a 9,50 €/kg mezclando café de 12 €/kg con café de 8 €/kg. ¿Cuántos kilos de cada tipo tenemos que poner?

\ex{3} Dos ciclistas salen a la vez de dos ciudades separadas 90 km y van el uno hacia el otro a 20 km/h y a 25 km/h. ¿Cuánto tardan en encontrarse? ¿A qué distancia de la primera ciudad?


% =====================================================
\seccio{6}{Ecuaciones de segundo grado incompletas}

\ex{1} Escribe cada ecuación en la forma general $ax^2+bx+c=0$, indica los coeficientes $a$, $b$ y $c$ y di si es completa o incompleta:\begin{tasks}(2)
  \task $5x-x^{2}=6$
  \task $3x^{2}=12$
  \task $x(x-4)=0$
  \task $(x+1)^{2}=2x+5$
\end{tasks}

\ex{1} Resuelve:\begin{tasks}(3)
  \task $x^{2}-36=0$
  \task $2x^{2}-50=0$
  \task $x^{2}+9=0$
  \task $4x^{2}-9=0$
  \task $3x^{2}=27$
  \task $x^{2}-7=0$
\end{tasks}\recorda{en $ax^2+c=0$ despejamos $x^2$ y hacemos la raíz: hay \textbf{dos} soluciones, $x=\pm\sqrt{\ }$, o ninguna si el número es negativo.}

\ex{1} Resuelve:\begin{tasks}(4)
  \task $x^{2}-5x=0$
  \task $3x^{2}+12x=0$
  \task $2x^{2}=7x$
  \task $-x^{2}+6x=0$
\end{tasks}\recorda{en $ax^2+bx=0$ sacamos factor común $x$: una solución siempre es $x=0$. ¡No dividas entre $x$!}

\ex{2} Reduce y resuelve:\begin{tasks}(2)
  \task $(x+3)(x-3)=7$
  \task $x(x-2)=4-2x$
  \task $(2x-1)^{2}=1$
  \task $(x+2)^{2}-4=3x$
\end{tasks}

\ex{2} Resuelve:\begin{tasks}(2)
  \task $\dfrac{x^{2}}{3}=12$
  \task $\dfrac{x^{2}}{4}-\dfrac{x}{2}=0$
  \task $\dfrac{x^{2}+1}{5}=2$
  \task $\dfrac{x^{2}-2}{2}-\dfrac{x^{2}+1}{3}=1$
\end{tasks}

\ex{3} Marta resuelve $5x^2=3x$ dividiendo los dos miembros entre $x$ y obtiene $5x=3$, es decir, $x=\tfrac35$. ¿Qué ha hecho mal? Resuélvela correctamente.

\ex{3} Escribe una ecuación de segundo grado incompleta que tenga por soluciones:\begin{tasks}(3)
  \task $4$ y $-4$
  \task $0$ y $3$
  \task $0$ y $-\tfrac12$
\end{tasks}


% =====================================================
\seccio{7}{Ecuaciones de segundo grado completas}

\ex{1} Resuelve con la fórmula:\begin{tasks}(4)
  \task $x^{2}-6x+8=0$
  \task $x^{2}+2x-8=0$
  \task $x^{2}-4x-5=0$
  \task $x^{2}+7x+12=0$
\end{tasks}\recorda{$x=\dfrac{-b\pm\sqrt{b^2-4ac}}{2a}$. \ Cuidado con los signos: si $b=-6$, entonces $-b=6$ y $b^2=36$.}

\ex{1} Resuelve:\begin{tasks}(4)
  \task $2x^{2}-5x+2=0$
  \task $3x^{2}+x-2=0$
  \task $6x^{2}-x-1=0$
  \task $4x^{2}+4x-3=0$
\end{tasks}

\ex{2} Resuelve (si $a$ es negativo, multiplica por $-1$; si todos los coeficientes tienen un divisor común, simplifica):\begin{tasks}(4)
  \task $-x^{2}+x+6=0$
  \task $2x^{2}+4x-6=0$
  \task $x^{2}-x-20=0$
  \task $3x^{2}-3x-18=0$
\end{tasks}

\ex{2} Reduce a la forma general y resuelve:\begin{tasks}(2)
  \task $(x-1)(x+4)=6$
  \task $(x+2)^{2}=3x+10$
  \task $x(x+3)=2(x+3)$
  \task $3x(x-2)-(x-1)(x+1)=9$
\end{tasks}

\ex{3} Resuelve:\begin{tasks}(2)
  \task $\dfrac{x^{2}}{2}-\dfrac{x}{3}=\dfrac{4}{3}$
  \task $\dfrac{x^{2}-1}{3}+\dfrac{x}{2}=2$
\end{tasks}

\ex{3} Resuelve:\begin{tasks}(2)
  \task $\dfrac{x(x-1)}{3}=x-1$
  \task $\dfrac{x(x+1)}{2}-\dfrac{(x-1)(x+2)}{3}=1$
\end{tasks}

\ex{3} Escribe una ecuación de segundo grado, en forma general, que tenga por soluciones:\begin{tasks}(2)
  \task $2$ y $-5$
  \task $\tfrac12$ y $3$
\end{tasks}\recorda{si las soluciones son $r$ y $s$, la ecuación $(x-r)(x-s)=0$ las tiene.}


% =====================================================
\seccio{8}{El discriminante}

\ex{1} Calcula el discriminante y di cuántas soluciones tiene cada ecuación, sin resolverla:\begin{tasks}(4)
  \task $x^{2}-10x+25=0$
  \task $2x^{2}+3x+4=0$
  \task $x^{2}-x-1=0$
  \task $-3x^{2}+2x+1=0$
\end{tasks}\recorda{$\Delta=b^2-4ac$. \ Si $\Delta>0$, dos soluciones; si $\Delta=0$, una (doble); si $\Delta<0$, ninguna.}

\ex{1} ¿Cuántas soluciones tiene cada ecuación?\begin{tasks}(3)
  \task $x^{2}+4x+4=0$
  \task $x^{2}+x-12=0$
  \task $5x^{2}-2x+1=0$
  \task $x^{2}-9=0$
  \task $4x^{2}-12x+9=0$
  \task $x^{2}+3x=0$
\end{tasks}

\ex{2} Resuelve (fíjate en el discriminante antes de continuar):\begin{tasks}(4)
  \task $9x^{2}-6x+1=0$
  \task $x^{2}+3x+5=0$
  \task $x^{2}-2x-2=0$
  \task $x^{2}+4x+1=0$
\end{tasks}

\ex{2} Escribe un valor de $c$ para que la ecuación $x^2+6x+c=0$ tenga: a) dos soluciones; b) una sola solución; c) ninguna solución.

\ex{3} Responde:\begin{tasks}(1)
  \task Halla $k$ para que la ecuación $x^2-4x+k=0$ tenga una única solución. ¿Cuál es?
  \task Halla $m$ para que $x^2+mx+16=0$ tenga una solución doble.
\end{tasks}

\ex{3} Responde:\begin{tasks}(1)
  \task ¿Para qué valores de $m$ la ecuación $x^2-2x+m=0$ tiene dos soluciones?
  \task Halla $k$ para que $x=3$ sea solución de $x^2+kx-6=0$ y calcula la otra solución.
\end{tasks}


% =====================================================
\seccio{9}{Problemas con ecuaciones de segundo grado}

\ex{1} Si al cuadrado de un número le restamos 5, obtenemos 44. ¿Qué número puede ser?\recorda{la ecuación puede tener dos soluciones: comprueba cuáles tienen sentido en el problema (una longitud o una edad no pueden ser negativas).}

\ex{1} El área de un cuadrado es 144 cm². ¿Cuánto mide el lado? ¿Y si el área fuera 50 cm²?

\ex{1} Un rectángulo mide 3 m más de largo que de ancho y su área es 40 m². ¿Qué dimensiones tiene?

\ex{2} Halla dos números naturales consecutivos cuyo producto sea 156.

\ex{2} Si al cuadrado de un número le restamos su triple, obtenemos 40. ¿Qué número puede ser?

\ex{2} Los catetos de un triángulo rectángulo miden $x$ y $x+7$ cm, y la hipotenusa 13 cm. ¿Cuánto miden los catetos?

\ex{2} Si al cuadrado de la edad de Pau le restamos el triple de su edad, obtenemos 130. ¿Cuántos años tiene?

\ex{2} Si aumentamos en 3 cm el lado de un cuadrado, su área aumenta 51 cm². ¿Cuánto mide el lado?

\ex{3} La suma de los cuadrados de dos números consecutivos es 85. ¿Cuáles pueden ser?

\ex{3} Un rectángulo tiene 34 m de perímetro y 60 m² de área. ¿Qué dimensiones tiene?

\ex{3} La hipotenusa de un triángulo rectángulo mide 10 cm y un cateto mide 2 cm más que el otro. ¿Cuánto miden los catetos?

\ex{3} Una piscina rectangular de 10 m × 6 m está rodeada por una acera de anchura constante. El área total (piscina y acera) es 96 m². ¿Qué anchura tiene la acera?

\vfill
\begin{center}\color{muted}\small\faCheckCircle[regular]\ \ ¡Ánimo! Comprueba tus resultados con las herramientas interactivas de los apuntes de la unidad.\end{center}

\end{document}
